\subsection{UNIVERSAL SETS AND MAPPINGS}

Let \(\mathfrak{G}\) be a theory which is stronger than the theory of sets, and let \(\mathbf{E}\) be a term in \(\mathfrak{G}\) . Let \(\Sigma\) be a species of structures in \(\mathfrak{G}\) . For simplicity we shall suppose throughout that \(\Sigma\) is defined on a single (principal) base set, and for brevity we shall say `` \(\Sigma\) -set'' for ``set endowed with a structure of species \(\Sigma\) ''. Furthermore, we shall suppose that the \(\sigma\) -morphisms have been defined for the species \(\Sigma\) ( \(\S 2\) , no. 1; as in \(\S 2\) , we shall say ``morphism'' in place of `` \(\sigma\) -morphism''). Finally, the species \(\Sigma\) being defined on the base set \(x\) and having \(s\) as generic structure ( \(\S 1\) , no. 4), let us suppose that a term \(\alpha \{x, s\}\) is defined in \(\mathfrak{G}_{\Sigma}\) , satisfying the following conditions:

\((\mathbf{Q}\mathbf{M}_{\mathbf{I}})\) The relation \(\alpha \{x,s\} \subset \mathcal{F}(\mathbf{E};x)\) is true in \(\mathfrak{C}_{\Sigma}\) .

(QM \(_{II}\) ) If (in a theory \(\mathfrak{G}'\) which is stronger than \(\mathfrak{G}\) ) F and F' are two sets endowed with structures \(\mathscr{S}\) , \(\mathscr{G}'\) of species \(\Sigma\) , and if f is a morphism of F into F', then the relation \(\varphi \in \alpha \{F, \mathscr{G}\}\) implies \(f \circ \varphi \in \alpha \{F', \mathscr{G}'\}\) .

We shall express the relation \(\varphi\in\alpha\{x,s\}\) by saying that \(\varphi\) is an \(\alpha\) -mapping of E into \(x\) (endowed with \(s\) ).

A \(\Sigma\) -set \(\mathbf{F}_{\mathbf{E}}\) and an \(\alpha\) -mapping \(\varphi_{\mathbf{E}}\) of \(\mathbf{E}\) into \(\mathbf{F}_{\mathbf{E}}\) are said to be universal if the following condition is satisfied:

(AU) For each \(\alpha\) -mapping \(\varphi\) of \(\mathbf{E}\) into a \(\Sigma\) -set \(\mathbf{F}\) there exists a unique morphism \(f\) of \(\mathbf{F}_{\mathbf{E}}\) into \(\mathbf{F}\) such that \(\varphi = f \circ \varphi_{\mathbf{E}}\) .

The pair \((\mathbf{F}_{\mathbf{E}}, \varphi_{\mathbf{E}})\) is then also said to be a solution of the universal mapping problem for E (relative to \(\Sigma\) , \(\sigma\) , and \(\alpha\) ).

Let \((\mathbf{F}_{\mathbf{E}}^{\prime},\varphi_{\mathbf{E}}^{\prime})\) and \((\mathbf{F}_{\mathbf{E}}^{\prime \prime},\varphi_{\mathbf{E}}^{\prime \prime})\) be two solutions of the universal mapping problem for E. The condition (AU) shows then that there exists a unique morphism \(f_{1}\) of \(\mathbf{F}_{\mathbf{E}}^{\prime}\) into \(\mathbf{F}_{\mathbf{E}}^{\prime \prime}\) and a unique morphism \(f_{2}\) of \(\mathbf{F}_{\mathbf{E}}^{\prime \prime}\) into \(\mathbf{F}_{\mathbf{E}}^{\prime}\) such that \(\varphi_{\mathbf{E}}^{\prime \prime} = f_{1}\circ \varphi_{\mathbf{E}}^{\prime}\) and \(\varphi_{\mathbf{E}}^{\prime} = f_{2}\circ \varphi_{\mathbf{E}}^{\prime \prime}\) . We have therefore \(\varphi_{\mathbf{E}}^{\prime} = f_{2}\circ f_{1}\circ \varphi_{\mathbf{E}}^{\prime}\) and \(\varphi_{\mathbf{E}}^{\prime \prime} = f_{1}\circ f_{2}\circ \varphi_{\mathbf{E}}^{\prime \prime}\) . Applying (AU) to the case where \(\mathbf{F} = \mathbf{F}_{\mathbf{E}}^{\prime}\) and \(\varphi = \varphi_{\mathbf{E}}^{\prime}\) , we find that \(f_{2}\circ f_{1}\) is the identity mapping of \(\mathbf{F}_{\mathbf{E}}^{\prime}\) onto itself. Similarly, \(f_{1}\circ f_{2}\) is the identity mapping of \(\mathbf{F}_{\mathbf{E}}^{\prime \prime}\) onto itself. Consequently (§ 2, no. 1, criterion CST8) \(f_{1}\) is an isomorphism of \(\mathbf{F}_{\mathbf{E}}^{\prime}\) onto \(\mathbf{F}_{\mathbf{E}}^{\prime \prime}\) , and \(f_{2}\) is its inverse isomorphism. This result is expressed by saying that the solution of the universal mapping problem for E is unique up to isomorphism.

To verify that a pair \((\mathbf{F}_{\mathbf{E}}, \varphi_{\mathbf{E}})\) is a solution of the universal mapping problem for E, it is often convenient to verify the following two conditions:

(AU \(_{I}'\) ) For every Σ-set F and every α-mapping φ of E into F, there exists a morphism f of F \(_{E}\) into F such that φ = f∘φ \(_{E}\) .

(AU \(_{II}'\) ) For every Σ-set F, two morphisms of F \(_{E}\) into F which agree on φ \(_{E}\) (E) are equal.

For if these two conditions are satisfied, the morphism \(f\) whose existence is ensured by \((\mathbf{AU}_{\mathbf{I}}^{\prime})\) is unique by \((\mathbf{AU}_{\mathbf{II}}^{\prime})\) . Conversely, it is clear that (AU) implies \((\mathbf{AU}_{\mathbf{I}}^{\prime})\) ; furthermore, if \(f\) and \(f'\) are two morphisms of \(\mathbf{F}_{\mathbf{E}}\) into \(\mathbf{F}\) which agree on \(\varphi_{\mathbf{E}}(\mathbf{E})\) , we have \(f \circ \varphi_{\mathbf{E}} = f' \circ \varphi_{\mathbf{E}}\) , whence \(f = f'\) by applying (AU) to the \(\alpha\) -mapping \(f \circ \varphi_{\mathbf{E}}\) . Hence (AU) implies \((\mathbf{AU}_{\mathbf{II}}^{\prime})\) .

